% Appendix F: the model answers to Paper A, from exam/paper_a_solutions.md.

% The answers' headings are labelled ansa:q:N.
\renewcommand{\answerprefix}{ansa}

\chapter{Model Answers to Paper A}
\label{ansa:ch:answers}

Marks are shown per part. Method carries them: a correct derivation with a slip in it is worth more
than a correct answer with no working, and later parts consume earlier ones, so an error should be
followed through rather than penalised twice. A candidate who decodes the register map wrongly in
Question~1 and then uses their own map consistently in Questions~4 and~6 loses the marks once.

Where a part asks for VHDL, mark the \textbf{hardware described}, not the syntax. Where it asks for
C++, mark the \textbf{byte order, the ordering of register accesses, and the qualifiers that carry
meaning}, not the punctuation.

Two parts, 2(d) and 4(d), ask what a provided testbench \emph{cannot} catch. An answer that treats a
passing testbench as proof has missed the point of the question however correct the rest of it is,
and those parts are marked accordingly.

% ==========================================================================================
\answersection{1}{The map, the top, and what a clean analysis does not prove}{12}

\answerpart{a}{3}
\code{0x0000000D} is \code{1101} in its low nibble.

\begin{booktable}{@{}lp{7.2em}lLl@{}}
  \thead{Bit} & \thead{Constant} & \thead{Position} & \thead{Meaning} & \thead{In \code{0x0D}} \\
  \midrule
  0 & \code{ST_TX_READY} & 0 & the TX FIFO is not full & \textbf{set} \\
  1 & \code{ST_RX_VALID} & 1 & the RX FIFO is not empty & clear \\
  2 & \code{ST_ERROR} & 2 & one or more \code{ERROR_FLAGS} bits set & \textbf{set} \\
  3 & \code{ST_TX_IDLE} & 3 & the TX FIFO empty \textbf{and} the line idle & \textbf{set} \\
\end{booktable}
\credit{1 mark for all four names and positions, correctly read against the value}

\textbf{The state.} The transmitter has nothing to send and nothing in flight --- TX-idle implies
TX-ready, which is why both are set --- no received byte is waiting, and at least one error has been
latched since the flags were last cleared. In this book's build that error can only be a framing
error, since \code{ER_FRAMING} is the only flag with a producer.

\textbf{What each call does.}
\begin{itemize}
  \item \code{write(0x41)} reads \code{STATUS}, finds \code{TX_READY} set, writes \code{0x41} to
    \code{TX_DATA} and returns \code{true}. Two transactions, ten bytes.
  \item \code{read(byte)} reads \code{STATUS}, finds \code{RX_VALID} clear, and returns
    \code{false}. It performs \textbf{no} \code{RX_DATA} read and \textbf{no} \code{RX_POP} write
    --- that omission is the contract, not an optimization, and it is what one of the host tests
    exists to check.
\end{itemize}
\worth{1}

\textbf{\code{ERROR_FLAGS}.} Index \textbf{6}, offset \textbf{\code{0x18}} (the index is the offset
divided by four). The command byte to read it is \textbf{\code{0x06}}; to write it,
\textbf{\code{0x86}}, bit~7 being the write bit. \worth{1}

\answerpart{b}{3}
\textbf{\code{ghdl -a} accepts it.} Both ports are \code{in std_logic}, so swapping them produces a
perfectly legal entity: the analyzer checks that each formal has a type and a direction, and both
still do. There is nothing here for it to object to. \worth{1}

\textbf{Elaboration accepts it too.} \code{uart_top_tb} binds positionally, so it connects its own
\code{sclk} signal to the entity's third port and its \code{mosi} signal to the fourth. Names never
enter into it. Every type and direction matches, so elaboration succeeds --- and the same is true
inside \code{uart_top}, where the architecture's positional \code{port map} into \code{spi_slave}
now hands the SPI clock line to \code{mosi} and the data line to \code{sclk}.

\textbf{Where it surfaces.} The first check anywhere in the course that fails is
\textbf{\code{uart_top_tb}, at the end of \chapref{5}} --- the first time the system testbench runs
at all. Before that it is \emph{skipped}, not passed: its datapath blocks do not exist. So a wiring
mistake made in \chapref{1} sits undetected for four chapters and then arrives as an SPI transaction
that produces nothing. \worth{1}

\textbf{The change that would move it forward.} Bind by \textbf{name} rather than by position, in
the testbench and in every instantiation. With named association the entity's port \emph{order}
stops mattering: both the testbench and \code{uart_top}'s architecture refer to \code{sclk} by
name, so moving it from position three to position four changes nothing and the transposition
becomes a no-op instead of a bug.

The cost is the reason the book did not: named association is verbose, it lets the declared order
drift from the documented one, and it changes what ``the contract'' means. This book made the
\textbf{order and the types} the contract precisely so that a port table in a chapter is a complete
specification of a module --- rename every signal and everything still binds. That is a real benefit,
and this failure is its real price. \worth{1}

Accept ``give the ports distinct subtypes so that a swap is a type error'' as a second answer, with
the observation that it is impractical for a bus of \code{std_logic} pins.

\answerpart{c}{4}
\noindent\textbf{(i)}\enspace\textbf{The mechanism is driver resolution.} \code{std_logic} is a
\emph{resolved} type: a signal with more than one driver takes the value its resolution function
computes from all of them, bit by bit. The placeholder is a concurrent assignment and therefore a
driver; \code{uart_regs}' \code{reg_rdata} output is a second one. Resolving \code{'1'} against
\code{'0'} gives \code{'X'}; resolving \code{'0'} against \code{'0'} gives \code{'0'}. For a
read-back of \code{0x00000004}, the low four bits carry:

\begin{narrowtable}{@{}cccc@{}}
  \thead{Bit} & \thead{Bank drives} & \thead{Placeholder drives} & \thead{Resolved} \\
  \midrule
  3 & \code{'0'} & \code{'0'} & \code{'0'} \\
  2 & \code{'1'} & \code{'0'} & \textbf{\code{'X'}} \\
  1 & \code{'0'} & \code{'0'} & \code{'0'} \\
  0 & \code{'0'} & \code{'0'} & \code{'0'} \\
\end{narrowtable}
\worth{1}

\textbf{What the testbench reports.} The \code{'X'} propagates out through the bridge and
\code{spi_slave} onto \code{MISO} and into the testbench's captured word, so
\code{rd(15 downto 0) = x"0004"} is false and the \code{BAUD_DIV} read-back assertion fires. The
symptom generalizes: \textbf{every bit the bank drives \code{'1'} comes back unknown}, so the value
looks like the one you wrote with all its ones missing. That is the signature to recognize, and the
fix is to delete the \chapref{1} placeholder line, exactly as Exercise~\ref{c5:ex:2} instructs.
\worth{1}

An answer that says ``two drivers, so it fails'' without naming resolution or without the per-bit
result gets one of the two marks.

\medskip\noindent\textbf{(ii)}\enspace\textbf{The value and the message.} \code{to_integer} on an
all-\code{'U'} vector prints

\begin{console}
NUMERIC_STD.TO_INTEGER: metavalue detected, returning 0
\end{console}

\noindent and returns \textbf{0}. \worth{1}

\textbf{Why 0 is worse than wrong.} \code{baud_gen}'s \code{div} port is declared
\code{natural range 1 to 65535}, so 0 is not a bad baud rate --- it is \textbf{outside the port's
subtype}, and binding it is a run-time bound-check failure that aborts the simulation. A wrong
divider you could at least see on a waveform; this one stops the run before there is a waveform. And
the library chose the value, not the designer, which is the deeper objection: a warning that
silently substitutes a value has made a design decision on your behalf.

\textbf{The guard.}

\begin{vhdlcode}
baud_div_int <= 1 when is_x(baud_div) or unsigned(baud_div) = 0
              else to_integer(unsigned(baud_div));
\end{vhdlcode}

\begin{itemize}
  \item \code{is_x(baud_div)} covers the \textbf{metavalue} case: before \chapref{5} nothing drives
    \code{baud_div}, so it sits at all-\code{'U'}.
  \item \code{unsigned(baud_div) = 0} covers the \textbf{legitimate zero} that arrives once
    \code{uart_regs} exists: \code{BAUD_DIV} holds whatever the bank reset it to until software
    writes a divider, and that reset value is a real, well-defined \code{0}.
\end{itemize}
\worth{1}

\code{1} is chosen because it is the smallest value the port accepts. Nothing useful is transmitted
at either value, so the only thing that matters is that the design elaborates and runs. Worth a
comment rather than a mark: \code{or} on \code{BOOLEAN} is short-circuit in VHDL, so the comparison
is never evaluated on the all-\code{'U'} vector, and the metavalue case stays warning-free as well
as legal.

\answerpart{d}{2}
\textbf{The port is \code{sync}'s \code{reset_s2_n} \emph{input}.} \code{sync} \textbf{consumes} the
synchronized reset; a module that requires a clean reset as an input cannot be the module that
manufactures one. Try it and the dependency is circular: \code{sync}'s own two flops would need
\code{reset_s2_n} to start from a known state, and \code{reset_s2_n} is the thing they were
supposed to produce. \worth{1}

\textbf{What the pattern buys.}
\begin{itemize}
  \item \textbf{Assert asynchronously.} The instant \code{reset_n} falls, every flop clears with no
    clock edge needed. That matters because the clock may not be running, or may be about to start,
    and a reset that requires a clock cannot rescue a design whose clock has not arrived.
  \item \textbf{Release synchronously.} The release is lined up with \code{clock} by
    \code{reset_sync}'s own two flops, so every flip-flop in the design leaves reset on the same
    edge.
\end{itemize}

Released asynchronously, the rising edge of \code{reset_n} could land inside some flops'
setup-and-hold window and not others'. Those flops would leave reset one cycle later than the rest,
so the design starts \textbf{half-reset}: a state machine already in \code{STATE_IDLE} beside a tick
counter that never cleared, or a FIFO whose \code{head} reset while its \code{tail} did not. That is
a state the design's own logic never anticipates, and it happens at a rate that depends on where the
button bounced. \worth{1}

% ==========================================================================================
\answersection{2}{The wire, the divider, and the transmitter}{13}

\answerpart{a}{4}
\code{0x53} is \code{0101 0011}, so \code{data(7 downto 0) = "01010011"}.

\textbf{The frame.} Concatenation puts the leftmost item in the highest bits, so
\code{frame <= STOP_BIT & data & START_BIT} gives

\begin{console}
frame(9 downto 0) = '1' & "01010011" & '0' = "1010100110"
\end{console}

\begin{booktable}{@{}lcccccccccc@{}}
  index & 9 & 8 & 7 & 6 & 5 & 4 & 3 & 2 & 1 & 0 \\
  \midrule
  value & 1 & 0 & 1 & 0 & 1 & 0 & 0 & 1 & 1 & 0 \\
  is & stop & \code{d7} & \code{d6} & \code{d5} & \code{d4} & \code{d3} & \code{d2} & \code{d1} &
    \code{d0} & start \\
\end{booktable}
\worth{1}

\textbf{The wire order.} \code{bit_idx} walks 0 upward, so the line carries \code{frame(0)} first:

\begin{booktable}{@{}lcccccccccc@{}}
  time order & 1 & 2 & 3 & 4 & 5 & 6 & 7 & 8 & 9 & 10 \\
  \midrule
  \code{frame} index & 0 & 1 & 2 & 3 & 4 & 5 & 6 & 7 & 8 & 9 \\
  level & 0 & 1 & 1 & 0 & 0 & 1 & 0 & 1 & 0 & 1 \\
  is & start & \code{d0} & \code{d1} & \code{d2} & \code{d3} & \code{d4} & \code{d5} & \code{d6} &
    \code{d7} & stop \\
\end{booktable}
\credit{2 marks: 1 for the order, 1 for the levels}

That is the whole point of building the frame as a vector: read it from index 0 upward and it
\emph{is} the wire order, start bit first and data least significant first, so the transmitter never
thinks about bit order again.

\textbf{The timing.} Ten bits at sixteen ticks each is \textbf{160 \code{baud_tick}s}. At
115200~baud \code{BAUD_DIV = 27}, so that is $160 \times 27 = $ \textbf{4320 clock cycles}, or
$4320 \times 20$\,ns $=$ \textbf{86.4\,µs}. \worth{1}

\answerpart{b}{3}
\begin{console}
exact divider = 50_000_000 / (16 x 57_600) = 50_000_000 / 921_600 = 54.2535
BAUD_DIV      = round(54.2535) = 54
\end{console}
\worth{1}

\begin{console}
tick rate  = 50_000_000 / 54 = 925_925.9 Hz
baud rate  = 925_925.9 / 16  = 57_870.4 baud
error      = (57_870.4 - 57_600) / 57_600 = +0.47%   (fast)
\end{console}
\worth{1}

\textbf{Why it is tolerated, and what has to be true.} The receiver does not run free: it hunts for
the falling edge of every start bit and restarts its own tick counter there. So the error does not
accumulate across a stream of bytes, only across the ten bit periods of a single frame.

The condition is therefore specific: the \textbf{stop-bit sample must still land inside the stop
bit}, nine and a half bit periods after the edge that started the frame. That is the last and
worst-placed sample in the frame, so it is the one that decides. Question~3 puts the budget at
roughly plus or minus five percent combined for an ideally centred sample; 0.47\,\% at this end
leaves the peer almost all of it. \worth{1}

An answer that says only ``it resynchronizes each frame'' earns half; the mark is for naming the
stop bit as the condition.

\answerpart{c}{3}
\textbf{When \code{tx} goes low.} \code{start} is looked at only in \code{STATE_IDLE}, and
\code{tx <= '0'} is scheduled in the same clocked process at the same rising edge, so \code{tx} is
low from just after \textbf{the very edge on which \code{start} was seen}. Nothing waits for a
\code{baud_tick}. The start bit therefore begins on a \emph{clock} edge, at whatever phase
\code{baud_gen}'s internal counter happens to be sitting. \worth{1}

\textbf{How long the start bit lasts.} The bit ends on the sixteenth \code{baud_tick} after entry to
\code{STATE_SEND}: tick~1 finds \code{ticks = 0} and increments, \dots, tick~15 finds
\code{ticks = 14} and increments to 15, and tick~16 finds \code{ticks = 15}, satisfies
\code{ticks >= MAX_TICKS - 1}, and advances \code{bit_idx}.

The \textbf{first} of those sixteen ticks arrives between 1 and \code{div} clocks after entry,
because the divider's counter was already partway through a window. Every tick after it is exactly
\code{div} clocks apart. One clock more follows: \code{tx <= frame(bit_idx)} reads \code{bit_idx}
before the edge that advances it, so the line changes to data bit~0 one clock after the sixteenth
tick. So, for $k$ in $1 \ldots \mathtt{div}$:

\begin{console}
start bit length = k + 15 x div + 1  clocks
                 = 15 x div + 2       at shortest
                 = 16 x div + 1       at longest
\end{console}
\worth{1}

A nominal bit is \code{16 x div} clocks, so the start bit is at worst \code{div - 2} clocks short,
or one clock long: under one tick period either way, that is under \textbf{one sixteenth of a bit}.
The spread of \code{div - 1} clocks is the phase of \code{baud_gen}'s counter when \code{start}
arrives. Simulated at \code{div = 4}, the start bit measures 62 to 65 clocks and every data bit
exactly 64. Accept \code{15 x div + 1} to \code{16 x div} from a candidate who misses the one-clock
lag but explains the tick phase; the variation is the point.

\textbf{Why nothing else inherits it.} \code{ticks} is cleared \emph{on} the tick that ends the bit,
and that tick is itself a tick boundary. So every bit after the first begins on a tick (one clock
after it, the same lag every time) and lasts exactly sixteen tick periods. The error is confined to
the start bit and does not accumulate down the frame.

And the receiver does not care: it triggers on the falling edge that opens the start bit and samples
the start bit half a bit later. A start bit up to one sixteenth short is still comfortably low at its
own midpoint, and every bit the receiver actually reads a value out of is full length. \worth{1}

\answerpart{d}{3}
\textbf{The two orderings.} \code{0xA5 = 1010 0101}, so $d_7 = 1$, $d_6 = 0$, $d_5 = 1$, $d_4 = 0$,
$d_3 = 0$, $d_2 = 1$, $d_1 = 0$, $d_0 = 1$.

\begin{booktable}{@{}Lcccccccc@{}}
  \thead{Order} & 1st & 2nd & 3rd & 4th & 5th & 6th & 7th & 8th \\
  \midrule
  least significant first ($d_0 \ldots d_7$) & 1 & 0 & 1 & 0 & 0 & 1 & 0 & 1 \\
  most significant first ($d_7 \ldots d_0$) & 1 & 0 & 1 & 0 & 0 & 1 & 0 & 1 \\
\end{booktable}

\textbf{They are identical.} A transmitter that packed its frame most significant first would put
exactly the same eight levels on the line, in exactly the same order, and \code{uart_tx_tb} would
have passed. \worth{1}

\textbf{The property.} \code{0xA5} is a \textbf{bit-reversal palindrome}: reverse its eight bits and
you get \code{0xA5} again. The constant chosen was one of the few bytes that cannot possibly
distinguish the two orderings.

\textbf{How many bytes share it.} A palindrome is determined entirely by its top four bits ---
$b_7$ fixes $b_0$, $b_6$ fixes $b_1$, $b_5$ fixes $b_2$, $b_4$ fixes $b_3$ --- so there are
$2^4 = $ \textbf{16} of them out of 256. Pick a byte at random and you have a 15-in-16 chance of
catching the bug; the earlier testbench had picked one of the 16. \worth{1}

\textbf{The constant the testbench uses today.} \code{0x53 = 0101 0011}, and its reverse is
\code{0xCA}:

\begin{booktable}{@{}Lcccccccc@{}}
  \thead{Order} & 1st & 2nd & 3rd & 4th & 5th & 6th & 7th & 8th \\
  \midrule
  least significant first ($d_0 \ldots d_7$) & 1 & 1 & 0 & 0 & 1 & 0 & 1 & 0 \\
  most significant first ($d_7 \ldots d_0$) & 0 & 1 & 0 & 1 & 0 & 0 & 1 & 1 \\
\end{booktable}

They differ at the \textbf{first} data bit, so the testbench fails immediately and reports
\textbf{data bit~0} mismatched. Any non-palindromic byte would do --- \code{0x01}, \code{0x0F},
\code{0x48}, \code{0xB2} --- and the mark is for seeing that asymmetry is the property required, not
for naming this particular byte.

\textbf{The lesson.} A green testbench proves the cases it contains, not the property you believed
you were checking. The check itself was always real --- the levels must match \code{TXBYTE} bit for
bit in the order expected --- but with a symmetric constant it had no power to distinguish an
ordering at all, so a genuine check had been quietly turned into a decoration by its own test data.
\worth{1}

Worth saying out loud, because the same reasoning applies to the other testbenches:
\code{uart_rx_tb} used \code{0xA5} and \code{0x3C}, and \code{uart_top_tb} used \code{0x5A}. All
three were palindromes, so at that revision \textbf{no testbench in \file{hw/} pinned down bit order
at all}. All four constants are non-palindromic today. One case survives even so:
\code{uart_top_tb} loops \code{tx} back to \code{rx}, so a transmitter and a receiver that are
reversed \textbf{together} cancel exactly and pass whatever constant it sends. Only the two unit
testbenches can catch that pair, which is exactly why a system-level loopback is not a substitute for
them.

Full marks for reaching that conclusion by any route. A candidate who writes out the two orderings,
sees that they match and says so has done the hard part.

% ==========================================================================================
\answersection{3}{The asynchronous line}{14}

\answerpart{a}{3}
\textbf{What can happen.} \code{rx} is governed by another chip, so it can change at any instant
relative to \code{clock}, including inside a flip-flop's \textbf{setup-and-hold window}. A flop
clocked in that window can go \textbf{metastable}: its output sits between the two logic levels for
an unbounded time before resolving, and resolves either way.

\textbf{What the failure actually is.} Sampling a signal mid-transition and getting either level
back is \textbf{not} the failure. Either answer is legitimate for a signal that genuinely was
changing, and one clock later you would have got the other one anyway. The failure is that for a
while the output is \textbf{not a level at all}. Two gates reading the same flop can resolve it
differently in the same cycle, so a state machine can take two branches at once and land in an
encoding that does not exist, and a counter can advance by something other than one. It is a
\emph{late} answer, not a wrong one, and lateness is what breaks a synchronous design. \worth{1}

\textbf{Why two flops.} The first flop is the one exposed to the asynchronous edge and may go
metastable. The second gives it a \textbf{full clock period} to settle before anything downstream
sees it, and the probability of a metastable state surviving that long falls exponentially with the
time allowed --- which is what the mean-time-between-failures expression describes. One flop gives it
no settling time whatever: whatever it is doing at the next edge is what the design consumes, so a
single register is not a synchronizer at all, however much it looks like one in simulation. Nothing
in a simulator is ever metastable, which is why \code{sync_tb}'s ``exactly two edges'' check exists.
\worth{1}

\textbf{Where it lives.} \code{uart_top} instantiates \code{sync} on the \code{rx} pin.
\code{uart_rx} deliberately does \textbf{not} contain one: its input port is named
\textbf{\code{rx_s2}}, and the \code{_s2} suffix is the announcement --- this signal has already
been through a two-flop synchronizer, so treat it as an ordinary synchronous input and sample it
directly.

The principle is that the module which owns the pins owns the crossing. \code{uart_top} owns the
pins, so it owns the crossing; \code{uart_rx} is then a purely synchronous design, which is easier to
reason about and easier to test. \worth{1}

\answerpart{b}{3}
\code{0x53 = 0101 0011}, so $d_0 = 1$, $d_1 = 1$, $d_2 = 0$, $d_3 = 0$, $d_4 = 1$, $d_5 = 0$,
$d_6 = 1$, $d_7 = 0$. Sent least significant first, that is the arrival order, and
\code{frame(bit_idx) <= rx_s2} with \code{bit_idx} counting up stores each where it belongs:

\begin{booktable}{@{}Lcccccccc@{}}
  \thead{Arrival} & 1st & 2nd & 3rd & 4th & 5th & 6th & 7th & 8th \\
  \midrule
  level & 1 & 1 & 0 & 0 & 1 & 0 & 1 & 0 \\
  \code{bit_idx} & 0 & 1 & 2 & 3 & 4 & 5 & 6 & 7 \\
\end{booktable}
\worth{1}

After the eighth, \code{frame(7 downto 0) = "01010011"}, and on a good stop bit that is what
\code{data_out} carries: \textbf{\code{0x53}}, the byte that was sent. \worth{1}

\textbf{Counting down.} With \code{bit_idx} walking 7 down to 0, the first-arriving bit lands in
index 7 and the last in index 0:

\begin{console}
frame = d0 d1 d2 d3 d4 d5 d6 d7 = 1 1 0 0 1 0 1 0 = "11001010" = 0xCA
\end{console}

\noindent The delivered byte is the \textbf{bit reversal} of the byte sent. Note that it is not
garbage and not noise: it is a well-formed byte, every frame arrives cleanly, \code{valid} pulses,
no \code{frame_err} is raised, and everything downstream works perfectly on data that is silently
wrong. \worth{1}

\answerpart{c}{5}
\noindent\textbf{(i)}\enspace The centre of bit~$n$ is $8 + 16n$ ticks after the true falling
edge. Against that:

\begin{booktable}{@{}Lp{9em}cc@{}}
  \thead{Sample} & \thead{Ticks after the true edge} & \thead{Ideal} & \thead{Past centre} \\
  \midrule
  start-bit re-check & 1 (notice) + 9 = 10 & 8 & \textbf{2 ticks} \\
  first data bit & 10 + 17 = 27 & 24 & \textbf{3 ticks} \\
  every later data bit & previous + 16 & +16 & \textbf{3 ticks} \\
  stop bit & last data + 17 = 156 & 152 & \textbf{4 ticks} \\
\end{booktable}

The offset is \textbf{fixed rather than accumulating} because data bit to data bit is exactly
sixteen ticks. The extra ticks are spent once at each \emph{state change} --- idle to start, start
to data, data to stop --- not once per bit. So the phase error is paid twice in the whole frame and
then held constant. \worth{1}

\medskip\noindent\textbf{(ii)}\enspace The stop bit is the tenth bit of an 8N1 frame, so $n = 9$.

\begin{console}
ideal sample tick   = 8 + 16 x 9 = 152
stop bit begins at  = 16 x 9     = 144
stop bit ends at    = 16 x 10    = 160
\end{console}
\worth{1}

Write $T_r$ for the receiver's tick period and $T_t$ for the transmitter's. The receiver samples 152
of \emph{its own} ticks after the edge; the transmitter's bit boundaries are 144 and 160 of \emph{its
own}. The sample stays inside the stop bit when
\[
  144\,T_t < 152\,T_r < 160\,T_t
  \quad\Longleftrightarrow\quad
  \frac{144}{152} < \frac{T_r}{T_t} < \frac{160}{152}
  \quad\Longleftrightarrow\quad
  0.9474 < \frac{T_r}{T_t} < 1.0526,
\]
so the combined mismatch budget is \textbf{plus or minus 5.26\,\%}, symmetric, about five percent
either way. \worth{1}

\medskip\noindent\textbf{(iii)}\enspace As built, the stop bit is sampled four ticks late, at
\textbf{156}:
\[
  144\,T_t < 156\,T_r < 160\,T_t
  \quad\Longleftrightarrow\quad
  \frac{144}{156} < \frac{T_r}{T_t} < \frac{160}{156}
  \quad\Longleftrightarrow\quad
  0.9231 < \frac{T_r}{T_t} < 1.0256,
\]
so \textbf{$-$7.69\,\% / +2.56\,\%}. \worth{1}

\textbf{Which side was eaten into.} $T_r / T_t > 1$ means the receiver's ticks are \emph{longer}
than the transmitter's --- the receiver is slow, or equivalently \textbf{the transmitter is fast}.
That limit falls from 5.26\,\% to 2.56\,\%, a little under half. The other side gains: a transmitter
running slow now has 7.69\,\% instead of 5.26\,\%, which is no use to anybody, because a real link is
limited by its worst direction.

\textbf{When it bites.} A peer transmitter that runs fast relative to this receiver. This is not
hypothetical: each end rounds its own integer divider and runs from its own clock, so each is off by
its own amount, and what the receiver sees is the difference. At 115200 this peripheral, with
\code{BAUD_DIV = 27}, is itself \textbf{+0.47\,\% fast}, which happens to push that difference
towards the roomy side; a divider that rounds up instead, making this end slow (9600, at
$-0.15$\,\%, is one), or a peer whose own rate errs upward, spends the small one. The margin is
still ample for a well-configured link --- the point is that the design started with half the
trailing margin the ideal tick-8 sample would have had, for three tick-sized implementation details
nobody chose deliberately. \worth{1}

Accept a candidate who works in ``percent baud error'' rather than tick-period ratios, provided the
asymmetry and the roughly halved fast-side figure come out.

\answerpart{d}{3}
\textbf{The FIFO.}

\begin{vhdlcode}
    rdata <= entries(tail);
    empty <= empty_s;
    full  <= full_s;

    full_s  <= '1' when count = DEPTH else '0';
    empty_s <= '1' when count = 0 else '0';

    process(clock, reset_s2_n) is
    variable push, pop: std_logic;
    begin
        if (reset_s2_n = '0') then
            head  <= 0;
            tail  <= 0;
            count <= 0;
        elsif (rising_edge(clock)) then
            push := wr and (not full_s);
            pop  := rd and (not empty_s);
            if (push = '1') then
                entries(head) <= wdata;
                advance_ptr(head);
            end if;
            if (pop = '1') then
                advance_ptr(tail);
            end if;
            if ((push = '1') and (pop = '0')) then
                count <= count + 1;
            elsif ((push = '0') and (pop = '1')) then
                count <= count - 1;
            end if;
        end if;
    end process;
\end{vhdlcode}
\credit{2 marks}

Four things carry the marks, and each is a place a plausible answer goes wrong:
\begin{itemize}
  \item \textbf{The guards are folded into \code{push} and \code{pop} once}, as variables, rather
    than repeated at each use. Everything downstream then tests the \emph{guarded} event, so a write
    to a full FIFO cannot move \code{head} and a read from an empty one cannot move \code{tail}.
  \item \textbf{\code{count} is adjusted only when the two disagree.} A simultaneous push and pop
    leaves it exactly where it was; incrementing in one branch and decrementing in another,
    unconditionally, is the classic bug here, and it desynchronizes the flags from the contents
    within one cycle.
  \item \textbf{\code{rdata} is a concurrent assignment from \code{tail}, not something the process
    writes.} That is the look-then-advance contract \chapref{5} depends on: the front byte is visible
    with no \code{rd} at all, and \code{rd} only advances.
  \item \textbf{Both flags come from \code{count}}, not from comparing \code{head} with \code{tail}
    --- which cannot tell full from empty in a ring buffer without a spare slot or an extra bit.
\end{itemize}
Accept a registered \code{rdata}, or \code{full} and \code{empty} derived from pointers plus a wrap
bit, only if the candidate says what they have paid for it. Accept \code{advance_ptr} written out
inline.

\textbf{How many frames.} \code{DEPTH = 8}, so the RX path absorbs \textbf{eight} whole frames. The
ninth \code{valid} pulse arrives with \code{full_s} set, so \code{push} is \code{'0'} and the byte
is \textbf{simply dropped} --- silently: no flag, no pulse, nothing on the wire. The line that
decides it is \code{push := wr and (not full_s)}. The FIFO's contract is to preserve order and keep
its flags honest, not to tell anybody it threw something away, which is why the caller must watch
the flags.

\textbf{The reserved bit.} \code{ER_OVERRUN}, \textbf{\code{ERROR_FLAGS} bit~2}. In this book's
build nothing in the hardware ever sets it, because only \code{frame_err} has a producer, so it reads
\code{0} unless software writes a \code{1} there; \code{ER_PARITY} (bit~1) is in the same situation.

That is a \textbf{reservation, not an omission}: the positions are fixed in the protocol, in
\file{uart_def.vhd} and in \file{register_map.hpp} (\code{error::OVERRUN = 2} on both sides), so an
implementation that adds overrun detection later drops straight in without renumbering anything or
breaking a driver that was already written against the map. A map that defined only the bits it
currently implements would force exactly that renumbering, and every existing driver with it.
\worth{1}

% ==========================================================================================
\answersection{4}{Ports become registers}{13}

\answerpart{a}{3}
\textbf{The assignments.}

\begin{vhdlcode}
    status_reg(31 downto 4) <= (others => '0');
    status_reg(ST_TX_READY) <= not tx_full_s;
    status_reg(ST_RX_VALID) <= not rx_empty_s;
    status_reg(ST_ERROR)    <= err_flags(0) or err_flags(1) or err_flags(2);
    status_reg(ST_TX_IDLE)  <= tx_empty_s and (not tx_busy);
\end{vhdlcode}
\worth{1}

Five concurrent assignments, not a process: there is no state here to clock. Accept a comparison
with \code{"000"}, or a helper function, in place of the OR for \code{ST_ERROR}, and accept any
spelling of the FIFO flag signals, but \textbf{the bits must be indexed by their \code{uart_def}
names}, not by literals --- that is the whole point of the package, and a candidate who writes
\code{status_reg(1)} has thrown away the protection the shared map exists to give. The reserved line
is worth having: without it bits 31--4 are undriven, and a read returns \code{'U'} in simulation
rather than the zeros Part~2 of the protocol (\secref{proto:part2}) promises.

\begin{booktable}{@{}lp{6.4em}p{8.6em}L@{}}
  \thead{Bit} & \thead{Constant} & \thead{Expression} & \thead{Why it cannot be stored} \\
  \midrule
  0 & \code{ST_TX_READY} & \code{not tx_full} & It is a property of the FIFO's own count, changing
    on every push \emph{and} every pop; a stored copy is wrong the moment one update is missed. \\
  1 & \code{ST_RX_VALID} & \code{not rx_empty} & The same --- and note it cannot be the receiver's
    \code{valid}, which is one clock wide and long gone before software polls. \\
  2 & \code{ST_ERROR} & the OR of the \code{ERROR_FLAGS} latches & It is a \emph{view} of other
    state, not a fact of its own; storing it means two places to keep in step. \\
  3 & \code{ST_TX_IDLE} & \code{tx_empty and (not tx_busy)} & It combines a FIFO flag with a
    datapath level, and no single writer knows both. \\
\end{booktable}
\credit{1 mark for reasoning that generalizes}

The general statement is worth the second mark on its own: each bit reports a \textbf{state}, while
the things that change that state are \textbf{pulses and edges}. A register written by events is a
register that is wrong whenever an event is missed, whenever two arrive together, and immediately
after reset unless its reset value is hand-maintained. Deriving the bit from the state it reports
removes all three failure modes by construction.

\textbf{The value.} The TX FIFO holds 2 of 8, so \code{tx_full = '0'} and bit~0 is \textbf{1}. The
RX FIFO is empty, so bit~1 is \textbf{0}. No error is latched, so bit~2 is \textbf{0}.
\code{tx_empty} is false (two bytes waiting) \emph{and} \code{tx_busy} is true, so bit~3 is
\textbf{0}. Everything above bit~3 reads zero.

\begin{console}
STATUS = 0x00000001
\end{console}
\worth{1}

\answerpart{b}{4}
\textbf{The failing check.} A read has no strobe, so a popping read can only be built from the
address: \code{rx_pop} also asserted whenever \code{reg_idx = REG_RX_DATA}. That pops on the edge
after which \code{uart_regs_tb} samples, so the first check to fire is \textbf{Case~5's
\code{RX_DATA} read-back}, which asserts that \code{RX_DATA} returns the byte just pushed and
reports \code{uart_regs_tb: RX_DATA mismatch, got 0x00!} --- the byte is gone before the testbench
looks. A variant that pops as the address \emph{leaves} \code{RX_DATA} survives Case~5 and trips
Case~6 instead, ``\code{RX_DATA read must not pop!}'', which reads \code{STATUS} after a bare
\code{RX_DATA} read and asserts that RX-valid is still set. Either is the split being pinned down;
award the mark for either, with what it asserts. \worth{1}

\textbf{Why a read with a side effect is harder.} The transport's abort rule says that if \code{SS}
rises before the fifth byte the transaction is abandoned with \textbf{no side effects}. Writes get
that for free: \code{spi_reg_bridge} asserts \code{reg_write} only when a write transaction
completes, so a \code{reg_write} the bank sees is always a real commit --- which is precisely what
lets \code{TX_DATA} and \code{RX_POP} be one-cycle actions with no state machine of their own.

A read has no such strobe. \code{reg_rdata} is \textbf{combinational}: the bank presents a value for
whatever \code{reg_addr} currently carries, and the bridge latches it once at the end of the command
byte and then shifts it out. The bank therefore cannot distinguish a read that completed from one
abandoned after the address byte --- and it cannot even tell how many times the same address was
presented.

So the thing the bridge would have to tell the bank, and currently has no way to say, is
\textbf{``this read transaction completed''} --- a read-side commit strobe, the mirror of
\code{reg_write}. The bank would then have to hold the pop pending until that strobe arrived and
drop it on an abort. That is the entire write-side machinery, duplicated, to make a read do
something a separate write already does for free. Splitting \code{RX_DATA} from \code{RX_POP}
sidesteps the question rather than solving it, which is why \secref{c5:app:a} calls it the single
most important idea in the module. \credit{2 marks}

\textbf{The driver trace.} The RX FIFO holds \code{0x41} then \code{0x42}.

\begin{booktable}{@{}cp{5.6em}Lp{7.2em}p{5em}l@{}}
  \thead{Call} & \thead{\code{STATUS} poll} & \thead{\code{RX_DATA} read} &
    \thead{\code{RX_POP} write} & \thead{Returns} & \thead{FIFO after} \\
  \midrule
  1 & RX-valid set & returns \code{0x41}, \textbf{and advances to \code{0x42}} & advances again,
    discarding \code{0x42} & \code{0x41}, \code{true} & empty \\
  2 & RX-valid clear & not issued & not issued & \code{false} & empty \\
\end{booktable}

\code{0x42} was \textbf{popped without ever being read}. The FIFO advanced twice for one delivered
byte, so every second byte is lost --- and lost invisibly, since nothing anywhere reports a discarded
entry. A testbench that sends one byte at a time never sees it. \worth{1}

\answerpart{c}{3}
\textbf{\code{uart_regs_tb}.} It does not write \code{CTRL} until its eighth case, so during the
transmit case the enable bit holds the bank's reset value, \code{0}. The \code{TX_DATA} write is
therefore swallowed, the TX FIFO stays empty, and the case that checks the FIFO front reports
\code{uart_regs_tb: TX FIFO front should be 0x5A, got 0x00!}. The message names the right register,
so the cause is findable. \worth{1}

\textbf{\code{uart_top_tb}.} It never writes \code{CTRL} at all. The byte never enters the TX FIFO,
so nothing is transmitted, nothing loops back, nothing arrives in the RX FIFO, and the testbench
polls \code{STATUS} for RX-valid up to its poll limit and gives up.

That message is the \textbf{less useful} of the two because it is a timeout: it says the byte never
came back and nothing about why. It is equally consistent with a broken transmitter, a broken
receiver, a broken FIFO, a wrong \code{BAUD_DIV}, a missing feeder, or a bank waiting to be enabled
--- exactly the failure \secref{c5:app:a} warns about, ``a timeout that says nothing about the
cause''. \worth{1}

\textbf{The fact that settles it.} \code{uart_regs} has \textbf{no \code{ctrl} output}. Read the port
list: there is no port by which \code{CT_ENABLE}, the parity select, the stop-bit select or either IRQ
mask could reach \code{uart_tx}, \code{uart_rx} or \code{baud_gen}. If no wire leaves the module
carrying \code{CTRL}, then nothing in this build may be gated on it, and the question is closed by
the interface rather than by taste. \worth{1}

\textbf{Why the bits exist anyway.} The register map is the \textbf{shared contract} with
\file{register_map.hpp} and with the protocol, and its positions are fixed so that an implementation
which wires them later stays compatible with both halves as already written. The driver's
\code{configure()} writes the enable bit for the same reason: the two sides agree on the wire, and
the write simply has no effect on the hardware yet.

\answerpart{d}{3}
\textbf{Why exactly one byte.} \code{busy} is a concurrent decode of the state:

\begin{vhdlcode}
busy <= '1' when state = STATE_SEND else '0';
\end{vhdlcode}

\noindent so it is high in the same simulation cycle as the state change the load caused. The
moment the transmitter accepts the byte it is in \code{STATE_SEND}, \code{tx_busy} is high, and
\code{tx_load} has already gone low --- one clock of \code{tx_load}, one \code{start} pulse, one
\code{tx_pop}.

It stays low for the rest of the frame because \code{tx_busy} stays high for the rest of the frame;
that is what holds it. And it cannot be \textbf{zero} loads, because the instant the transmitter
returns to \code{STATE_IDLE} with a non-empty FIFO both terms are true again and the next byte goes.
\worth{1}

\textbf{The simplification.} They are right that \code{uart_tx} ignores \code{start} outside
\code{STATE_IDLE}, so the transmitter is safe. The bug is in the \emph{other} line: \code{tx_pop} is
the same signal.

With \code{tx_load <= not tx_empty}, \code{tx_load} is high \textbf{continuously} for as long as
the FIFO has anything in it. So the FIFO is popped \textbf{once per 50\,MHz clock cycle}. A FIFO
holding four bytes empties in four consecutive clocks, about 80\,ns, while the first byte's start bit
is still on the wire and has roughly 8.6\,µs left to run.

\textbf{One} of the four bytes reaches the wire. The other three survive about 60\,ns at most and are
discarded unsent, with no flag and nothing to see. \worth{1}

\textbf{Does \code{uart_top_tb} catch it?} \textbf{No.} The testbench writes exactly one byte to
\code{TX_DATA}, so the FIFO empties after a single pop either way and the two designs behave
identically. It would take a testbench that queues two bytes to see the difference --- which is the
same lesson as Question~2(d) in different clothes: a passing testbench proves the cases it contains.
\worth{1}

Worth noting for a candidate who spots it: this is also why the feeder could not be written before
\chapref{5}. \code{tx_empty} comes from the register bank and \code{tx_busy} from the transmitter,
so it needs a block from \chapref{2} and a block from \chapref{5} to exist at the same time.

% ==========================================================================================
\answersection{5}{The driver's contracts}{12}

\answerpart{a}{4}
\begin{cppcode}
if (status & (1U << status::RX_VALID)) { /* a byte is waiting */ }
\end{cppcode}

\begin{vhdlcode}
if (status(ST_RX_VALID) = '1') then -- the same bit, indexed by position
\end{vhdlcode}
\credit{1 mark for both}

\textbf{\code{if (status & status::RX_VALID)}.} \code{status::RX_VALID} is \code{1}, so the mask is
\code{0x1} --- which is \textbf{\code{TX_READY}'s} bit. The position has been used where a mask was
meant, and it happens to be a legal mask for the wrong bit.

Consequence: \code{read()} returns \code{true} whenever the transmitter has room, which after reset
it always does. It reads \code{RX_DATA} from an empty FIFO, so the bank hands back whatever the ring
buffer holds at \code{tail} --- reset data or a stale entry --- and the driver returns that byte to
the caller as if it had arrived on the wire. It then issues an \code{RX_POP} that the FIFO's
\code{not empty} guard silently drops. \code{app::EchoNode} over this driver emits a stream of
garbage bytes at the speed of its poll loop, having received nothing at all. \worth{1}

\textbf{\code{writeReg(reg::CTRL, ctrl::ENABLE);}.} \code{ctrl::ENABLE} is \code{0}, so this writes
\code{0x00000000} to \code{CTRL} --- clearing the register rather than setting bit~0. The correct
form is \code{writeReg(reg::CTRL, 1U << ctrl::ENABLE)}.

Consequence, and this is the interesting part: \textbf{in this book's build, nothing happens.}
\code{CTRL} gates nothing, so the peripheral transmits and receives exactly as before, and every
testbench and every rung of the bench passes. Only the host suite notices ---
\code{Uart.ConfigureWritesBaudThenEnable} expects the \code{CTRL} write to end in \code{0x01} and
gets \code{0x00} --- and without that one assertion the bug is a time bomb: it becomes a peripheral
that never enables on the first day somebody wires \code{CT_ENABLE} to something real, in a driver
that has looked correct for months. \worth{1}

\textbf{A mask stored where the map stores a position.} \code{RX_VALID} is now \code{2}, so
\code{1U << status::RX_VALID} shifts by two and tests \textbf{bit~2}, the \code{ERROR} bit.
\code{read()} returns \code{true} when an error is latched and \code{false} otherwise, so it never
delivers a byte and does deliver on a framing error.

Worse than the behaviour is where the bug lives: the two sides of the wire now \textbf{disagree
about the map}, since \file{uart_def.vhd} still has \code{ST_RX_VALID = 1}. Nothing on the host can
find it, because the host tests script the stub with values produced from this same header --- the
driver and its tests are wrong together and agree perfectly. It is a bench-only bug, and ``both
sides store positions, transcribed from the protocol rather than from each other'' is the rule that
exists to prevent exactly it. \worth{1}

\answerpart{b}{3}
\textbf{The three methods.} \code{begin()}, \code{transfer(uint8_t)} returning \code{uint8_t}, and
\code{end()}. Select, exchange one byte in full duplex, deselect --- the framing of Part~3 of the
protocol and nothing more. \worth{1}

\textbf{What it does not know.} It knows nothing about \textbf{registers} --- no indices, no command
byte, no five-byte shape, no byte order --- and nothing about the \textbf{UART} --- no baud, no FIFOs,
no status bits. It is a dumb byte pipe.

\textbf{What the byte level buys.} All the protocol knowledge stays in one place, the driver, so
there is exactly one implementation of the transaction shape to get right and exactly one to test.
And the layer below becomes something a stub can imitate \textbf{perfectly}, because ``return a byte
for a byte'' has no behaviour to model --- there is nothing for a fake to get subtly wrong.

A seam at \code{readReg} and \code{writeReg} would push the command byte and the byte order
\emph{below} the line. The stub would then have to reimplement them, and a test asserting that the
driver produced the right transaction would really be asserting that the stub and the driver made
the same mistake. The byte-level seam is what makes the test meaningful, not merely possible.
\worth{1}

\textbf{The two implementations, and what changes.} \code{driver::transport::Stub} (\chapref{8},
host, scripted) and \code{driver::transport::AvrSpi} (\chapref{9}, target, real). What changes above
the seam when they swap: \textbf{nothing at all.} The register map, \code{Uart}, the blocking helpers
and \code{EchoNode} are recompiled, not rewritten --- which is the entire return on the design work
done in \chapref{6}. \worth{1}

\answerpart{c}{3}
\begin{enumerate}
  \item \textbf{\code{#include <cstdint>} and \code{std::uint8_t}, \code{std::uint32_t}.} The AVR
    target is \emph{freestanding} and avr-libc ships \textbf{no C++ standard library at all}, so
    there is no \code{<cstdint>} and no \code{std} namespace to qualify anything with. Replace them
    with \code{<stdint.h>} and the bare \code{uint8_t} and \code{uint32_t}.
  \item \textbf{\code{namespace driver::uart \{ ... \}}.} The compact nested-namespace definition is
    C++17, and the reference toolchain does not accept \code{-std=c++17} as an option at all.
    Replace it with nested blocks: \code{namespace driver \{ namespace uart \{ ... \} \}}.
  \item \textbf{\code{inline constexpr} at namespace scope.} \code{inline} \emph{variables} are also
    C++17. Plain \code{constexpr} is the replacement and costs nothing: at namespace scope it already
    has internal linkage, so each translation unit gets its own compile-time copy of a constant that
    never occupies storage.
  \item \textbf{\code{[[nodiscard]]}.} Rejected by that toolchain, so it is omitted from everything
    that has to cross-compile. The host-only test suites keep it, because they never build for the
    target.
\end{enumerate}

The point in each case is that the toolchain is missing a \textbf{language level and a library}, not
merely being conservative. The reference is the AVR/GNU C++ compiler bundled with Microchip Studio,
several GCC releases behind a desktop g++, so ``it compiled on my Linux avr-g++'' is not evidence
that it will build for the target the book names.
\credit{2 marks: 1 for the four rejections, 1 for reasons that name the missing language level or
library rather than ``it is old''}

\textbf{The corrected header.}

\begin{cppcode}
#pragma once

#include <stdint.h>

namespace driver
{
namespace uart
{
constexpr uint8_t TxReady{0U};

class Interface
{
public:
    virtual ~Interface() noexcept = default;

    virtual uint32_t status() const noexcept = 0;
};
} // namespace uart
} // namespace driver
\end{cppcode}
\worth{1}

\code{#pragma once} is not in the original listing and is not one of the four defects, but a header
without an include guard is a defect of its own and the book's headers all carry one; award the mark
without it, and note it. \code{virtual} and \code{= 0} survive untouched, as do \code{noexcept},
\code{override}, \code{final} and \code{= default} --- the subset is narrower than modern C++, not a
different language.

\answerpart{d}{2}
\textbf{What each \code{bool} means.} \code{write(byte)} returns \code{true} if the byte was accepted
into the TX FIFO and \code{false} if the FIFO was full, in which case \textbf{nothing was written}
--- the driver does not queue, retry or buffer. \code{read(byte&)} returns \code{true} if a byte was
placed in the out-parameter, and \code{false} if none was waiting, in which case no \code{RX_DATA}
read and no \code{RX_POP} were issued. \worth{1}

\textbf{The bits behind them.} \code{STATUS} bit~0, \code{TX_READY} (\code{not tx_full}), makes the
write answer possible; bit~1, \code{RX_VALID} (\code{not rx_empty}), makes the read answer
possible. The \chapref{5} mechanism behind both is that \code{STATUS} is \textbf{computed from the
FIFO flags}, so each is a level that stays true until the condition changes --- a pollable fact
rather than an event that has to be caught.

\textbf{The blocking versions.} The free functions \code{writeBlocking} and \code{readBlocking} in
\file{driver/uart/blocking.hpp}, taking \code{Interface&} so that they work over any implementation
and dispatch to the concrete driver at run time. They are free functions rather than methods because
they add no state and no new contract --- putting them in the interface would oblige every
implementation, including every stub, to reimplement a one-line spin loop.

If the interface itself blocked, two things would be lost. Every operation would become
non-deterministic to test: a test would have to arrange for the stub to \emph{eventually} succeed,
and a mistake would hang rather than fail. And an application could no longer poll ---
\code{EchoNode} could not re-check its \code{stop} flag between reads, which is the exact property
its test depends on. \worth{1}

% ==========================================================================================
\answersection{6}{The five-byte transaction}{13}

\answerpart{a}{4}
\textbf{The command bytes.}

\begin{booktable}{@{}p{7.6em}ccp{6.4em}L@{}}
  \thead{Bytes} & \thead{Bit 7} & \thead{Index} & \thead{Register} & \thead{Direction} \\
  \midrule
  \code{82 00 00 00 1B} & 1 & 2 & \code{BAUD_DIV} & write \code{0x0000001B} = 27 \\
  \code{81 00 00 00 01} & 1 & 1 & \code{CTRL} & write \code{0x00000001} = \code{1U << ctrl::ENABLE} \\
  \code{00 00 00 00 00} & 0 & 0 & \code{STATUS} & read, four dummy bytes \\
  \code{04 00 00 00 00} & 0 & 4 & \code{RX_DATA} & read, four dummy bytes \\
  \code{85 00 00 00 01} & 1 & 5 & \code{RX_POP} & write \code{1} \\
\end{booktable}
\worth{1}

\textbf{Which belong together.}
\begin{itemize}
  \item Transactions 1 and 2 are one \textbf{\code{configure(27)}}: the baud divider first, the
    enable afterwards, in that order, because the peripheral should know its rate before it is
    switched on.
  \item Transactions 3, 4 and 5 are one successful \textbf{\code{read(byte)}}: poll \code{STATUS},
    pure read of \code{RX_DATA}, separate \code{RX_POP} write. Three register accesses, in that
    order, which is the \chapref{5} contract showing up on the other side of the wire.
\end{itemize}
\worth{1}

\textbf{What the driver did with the returned bytes.} In both reads, the reply that came back during
the \emph{command} byte is discarded --- it is meaningless, since the slave does not yet know which
register is being asked for. The four bytes returned during the four dummy exchanges are assembled
\textbf{most significant byte first}: the first into bits 31--24, then 23--16, then 15--8, and the
last into bits 7--0. In transaction 3 the driver then tested bit~1 (\code{RX_VALID}) of the
assembled word; in transaction 4 it kept the low byte and copied it into the caller's out-parameter.
\worth{1}

\textbf{What the call counts prove.} They prove each transaction was \textbf{framed exactly once},
and the two counts being equal means every \code{SS} that fell also rose --- five complete, balanced
transactions rather than one long one.

The byte log alone cannot distinguish this from a driver that called \code{begin()} once, pushed
all twenty-five bytes, and called \code{end()} at the very end. On a stub that is invisible: the
same bytes in the same order. On real hardware it is fatal --- \code{spi_reg_bridge} would see
\code{ss_active} assert once, decode the first five bytes as one \code{BAUD_DIV} write, and treat
the remaining twenty as a continuation of a transaction that was already over, so nothing after the
first would be decoded at all. Framing is not visible in a list of MOSI bytes, so it has to be
counted separately, which is why the suite asserts on \code{beginCalls()} and \code{endCalls()} in
eight different cases. \worth{1}

\answerpart{b}{3}
\noindent\textbf{(i)}\enspace\textbf{The bytes.} Least significant first, \code{configure(27)}'s
first transaction puts

\begin{console}
82 1B 00 00 00
\end{console}

\noindent on the wire. \worth{1}

\textbf{What the peripheral commits.} The peripheral is not consulted about byte order --- the
protocol says the first data byte is bits 31--24. So it assembles \code{0x1B000000} and commits that
to \code{BAUD_DIV}. \code{BAUD_DIV} is bits 15--0, so \code{baud_div} carries \code{x"0000"}.

\code{uart_top}'s guarded conversion then hits its \textbf{second} case,
\code{unsigned(baud_div) = 0}, and substitutes \code{1}. \code{baud_gen} divides the 50\,MHz clock by
one, so it ticks every clock:

\begin{console}
tick rate = 50_000_000 Hz
baud rate = 50_000_000 / 16 = 3_125_000 baud
\end{console}

\noindent \textbf{3.125\,Mbaud}, about 27 times the intended rate. The line runs, \code{tx} toggles,
\code{uart_tx_tb} and \code{uart_regs_tb} are unaffected because neither involves the driver --- and
nothing at the far end can read a byte. \worth{1}

\medskip\noindent\textbf{(ii)}\enspace The peripheral shifts \code{0x0000000B} out most significant
first, so the four bytes returned are \code{00 00 00 0B}. Assembled least significant first, the
first goes to bits 7--0 and the last to bits 31--24:

\begin{console}
driver holds 0x0B000000
\end{console}

\noindent\code{status & (1U << status::TX_READY)} is then \code{0}. \code{write()} returns
\code{false} on every call for ever, and \code{writeBlocking()} --- a bare
\code{while (!uart.write(byte)) {}} --- spins for ever. The firmware hangs on the first byte it tries
to send, with the peripheral perfectly healthy and reporting itself ready. \worth{1}

\answerpart{c}{3}
\textbf{No \code{RX_POP}.} \code{RX_DATA} is a pure read, so the front byte is never discarded.
RX-valid stays set, and every subsequent \code{read()} polls, finds it set, re-reads the same front
byte and hands it back. \code{app::EchoNode} echoes the \textbf{first character it ever receives},
over and over, as fast as the poll loop runs --- the terminal fills with a single repeated character
and nothing typed afterwards has any effect. It is a very recognizable bench symptom, which is a
small mercy. \worth{1}

\textbf{Pop before read.} The front byte is discarded before it is looked at, so the \code{RX_DATA}
read returns whatever the FIFO shows \emph{afterwards}: the next byte if one has already arrived, or
stale empty-FIFO data if not. Every byte is either skipped outright or delivered one position late,
and the stream is silently corrupted rather than stalled --- which is harder to diagnose than the
first mistake, not easier. \worth{1}

\textbf{The pop on empty.} The test is \textbf{\code{Uart.ReadWhenEmptyReturnsFalseAndDoesNotPop}}:
\code{read()}, on no data, must report failure and issue no \code{RX_POP}, so the record holds one
transaction, not three. What the stray pop actually does to \emph{this} peripheral: \code{fifo}
guards \code{rd} on \code{not empty}, so it is dropped and no queued byte is lost.

``Nothing bad happens'' is not a defence, for three reasons:
\begin{itemize}
  \item The guard is an implementation detail of \code{fifo}, \textbf{not a promise the register
    map makes}. The protocol says \code{RX_POP} discards the front byte; it does not say what
    happens when there is not one.
  \item It costs a five-byte SPI transaction on every idle poll, which on a poll loop with an empty
    FIFO is most of the bus traffic and most of the latency.
  \item A driver that depends on an unspecified guard breaks against the first compliant peripheral
    that does not have one, and it will break on the bench rather than on the host.
\end{itemize}
\worth{1}

\answerpart{d}{3}
\textbf{Why the three are not \code{const}.} \code{begin()}, \code{transfer()} and \code{end()}
drive \code{SS} and shift real bytes through real hardware. They are \textbf{actions with side
effects}, not observations, and marking them \code{const} would be a lie about what calling them
does --- to the reader first and to the optimizer second. \worth{1}

\textbf{What the cast does, and why it is not needed here.}
\code{const_cast<transport::Interface&>(myTransport)} yields a non-\code{const} reference the three
methods can be called on --- but it strips nothing, because the enclosing \code{const} never reached
the transport.

The compiler would accept the calls \textbf{without} it as \code{Uart} is written today, because
\code{myTransport} is a \textbf{reference member}: \code{const} does not propagate through a
reference, so inside a \code{const} method \code{myTransport} still names a
\code{transport::Interface&}. The cast is documentation of intent rather than a requirement.

It becomes genuinely required the moment the transport is held differently --- by value, or behind a
\code{const}-propagating handle --- because then the object's \code{const} reaches the member and the
calls stop compiling. Writing it now means that change is a one-word edit rather than a redesign.
\worth{1}

\textbf{The undefined case.} Casting away \code{const} is undefined behaviour when the object being
cast is \textbf{actually \code{const}} --- declared \code{const}, or a temporary bound to a
\code{const} reference --- and is then modified through the resulting reference. That is not what
happens here: \code{myTransport} refers to a real, non-\code{const} transport that the caller
constructed and owns (a \code{Stub} on the host, an \code{AvrSpi} on the target), so driving it is
defined. The one hard rule is to cast away \code{const} only on an object that is not genuinely
\code{const}, and this obeys it. \worth{1}

% ==========================================================================================
\answersection{7}{The real transport, and the real toolchain}{11}

\answerpart{a}{4}
\textbf{What each statement achieves.}
\begin{itemize}
  \item \code{DDRB |= ...} makes \code{SCK} (PB5), \code{MOSI} (PB3) and \code{SS} (PB2)
    \textbf{outputs}.
  \item \code{PORTB |= (1U << SS)} idles the chip select \textbf{high}, that is deasserted, before
    any transaction.
  \item \code{SPCR = ...} enables the peripheral, selects master, and sets the prescaler.
\end{itemize}
\worth{1}

\textbf{How \code{SPCR} delivers the contract, and the bits left clear.}

\begin{booktable}{@{}p{3.6em}p{3.4em}L@{}}
  \thead{Bit} & \thead{State} & \thead{What it gives} \\
  \midrule
  \code{SPE} & set & The SPI peripheral is enabled. \\
  \code{MSTR} & set & Master: the ATmega generates \code{SCK}, and the FPGA is the slave. \\
  \code{SPR0} & set & With \code{SPR1} clear and \code{SPI2X} clear, the prescaler is 16:
    16\,MHz / 16 = \textbf{1\,MHz}. \\
  \code{DORD} & \textbf{clear} & \textbf{MSB first}, in every byte, as the protocol requires. \\
  \code{CPOL} & \textbf{clear} & \code{SCK} \textbf{idles low}. \\
  \code{CPHA} & \textbf{clear} & Sample on the \textbf{leading} edge. With \code{CPOL} clear that is
    the rising edge. \\
  \code{SPIE} & \textbf{clear} & No SPI interrupt. The transport polls, which is why it needs neither
    \code{<avr/interrupt.h>} nor any ISR. \\
  \code{SPR1} & \textbf{clear} & Half of the prescaler selection above. \\
\end{booktable}

\code{CPOL = 0} together with \code{CPHA = 0} \textbf{is} SPI mode 0, which is what Part~3 of the
protocol asks for. So all three requirements --- mode, order and rate --- are met by bits that are
\emph{not} set, which is exactly why an answer that lists only the three bits present is incomplete.
\credit{2 marks}

\textbf{Why \code{MISO} appears nowhere.} It is an input at reset and must stay one, so the
constructor has nothing to do to it. The book's rule is that the constructor touches only what it
must, so that the destructor can put back exactly that and nothing else --- \code{MISO} was never
taken, so it is never given back. \credit{half of the last mark}

\textbf{Without the middle statement.} \code{PORTB} resets to zero, so the instant \code{DDRB} makes
PB2 an output the pin is actively driven \textbf{low}: the chip select is asserted from construction
onward, with the bus idle. \code{spi_slave} sees \code{ss} fall and loads its first byte, and the
bridge's byte counter waits at zero. The first transaction therefore still decodes, provided
\code{SCK} has not moved: \code{begin()}'s drive low is a no-op, the five bytes arrive as a command
and four data bytes, and \code{end()} raises \code{SS} for the first time, after which every
transaction is framed normally. That is what makes the omission dangerous rather than loud: nothing
on the bench and nothing in a MOSI byte log looks wrong, while the contract's idle-high chip select
has been broken. Any \code{SCK} activity between construction and the first \code{begin()} would be
counted as the start of a command byte and misalign that whole first transaction, because the only
boundary the bridge knows is \code{SS} rising, which ties back to 6(a).
\credit{half of the last mark}

\answerpart{b}{3}
\textbf{What it returns.} The byte left in \code{SPDR} from the \textbf{previous} exchange. The
receive buffer still holds the last completed transfer's result, because the transfer just started
has not shifted a single bit yet --- the write to \code{SPDR} starts the shift, it does not perform
it. \worth{1}

\textbf{How long it takes.} Eight \code{SCK} periods, one per bit. At 1\,MHz that is \textbf{8\,µs},
which at 16\,MHz is \textbf{128 ATmega clock cycles} --- an eternity in instructions, and the reason
the omission is both easy to make and impossible to get away with. \worth{1}

\textbf{What the spin waits on, and the side effect.} The spin waits on \textbf{\code{SPIF}} in
\code{SPSR}:

\begin{cppcode}
SPDR = byte;
while (0U == (SPSR & (1U << SPIF))) {}
return SPDR;
\end{cppcode}

\noindent Reading \code{SPSR} while \code{SPIF} is set and \emph{then} reading \code{SPDR}
\textbf{clears \code{SPIF}}, arming the next transfer. That is not incidental: without it the flag
would still be set at the top of the next \code{transfer()} and the spin would fall straight
through, so the correct code owes its correctness to the clearing sequence as much as to the wait.

\textbf{What the driver sees.} On the host mock, which models only the missing wait, every returned
byte is one exchange stale: the reply to the command byte is discarded anyway, so \code{readReg}
assembles the replies that belonged to the command exchange and the first three data exchanges --- a
\code{STATUS} word built from the tail of the previous transaction and the head of this one. On the
ATmega it is worse, because the transmit side is single-buffered: a write to \code{SPDR} while a
byte is still shifting sets \code{WCOL} and is discarded. The whole of \code{readReg} runs in about
the time one byte takes at 1\,MHz, so the dummy writes collide with the command byte and are lost,
\code{end()} raises \code{SS} with the transaction nowhere near five bytes, and the bridge abandons
it with no side effects --- a \code{writeReg} built on this \code{transfer} commits nothing either.
Every \code{SPDR} read returns whatever the receive buffer last completed, never the register's
bytes, which is the worst kind of bug to meet first on a bench. \worth{1}

\answerpart{c}{2}
\textbf{What the hardware does.} It clears \textbf{\code{MSTR}} in \code{SPCR} and sets \code{SPIF}.
The ATmega demotes itself to a \textbf{slave} on the spot, stops driving \code{SCK}, and abandons
the transfer in flight. This is why \code{SS}/PB2 must be an output in master mode even when
something else is being used as the chip select --- and here PB2 \emph{is} the chip select, so
driving it is the transport's job anyway. \worth{1}

\textbf{What the spin does, and the symptom.} The spin in flight ends \textbf{immediately}, because
the demotion set \code{SPIF}, and \code{transfer()} returns a byte that was never clocked in: that
call succeeds, wrongly. The next one is where it stops. The SPI is now a slave, so writing
\code{SPDR} only loads the shift register and waits for a master's \code{SCK} that nobody drives;
\code{SPIF} never sets, and the spin loops for ever.

On the bench: \code{SCK} stops moving on the analyzer, and the firmware freezes inside
\code{transfer()} one byte after the demotion. Nothing recovers --- the transport configured
\code{SPCR} once in its constructor and never revisits it, so \code{MSTR} stays clear until the chip
is reset. \worth{1}

\answerpart{d}{2}
\textbf{Two host constructs that do not survive.} Any two of:
\begin{itemize}
  \item \textbf{Anything from the standard library, \code{std::cout} above all.} There are no
    iostreams, so debug output goes out the ATmega's \textbf{own} hardware USART: set \code{UBRR0}
    from \code{F_CPU} and the target baud, enable the transmitter in \code{UCSR0B}, then poll
    \code{UCSR0A & (1U << UDRE0)} and write \code{UDR0}. (Not to be confused with the FPGA
    peripheral this book builds.)
  \item \textbf{Dynamic allocation}, \code{new} or \code{std::make_unique}. There is no
    \code{operator new}, so everything is static or automatic storage --- which is why every seam in
    this book is a reference injected into an object living in \code{main}'s frame.
  \item \textbf{Exceptions and RTTI.} No \code{throw}, no \code{dynamic_cast}, no \code{typeid};
    errors become return values, which is what \code{write()} returning \code{false} on a full FIFO
    already is.
\end{itemize}
\worth{1}

\textbf{\file{env.cpp}.}

\begin{cppcode}
void operator delete(void*) noexcept {}
void operator delete(void*, unsigned int) noexcept {}

extern "C" void __cxa_pure_virtual() {}

extern "C" int  __cxa_guard_acquire(volatile void* g) { return !*(char*)g; }
extern "C" void __cxa_guard_release(volatile void* g) { *(char*)g = 1; }
extern "C" void __cxa_guard_abort(volatile void*) {}
\end{cppcode}

\noindent What emits each reference:
\begin{itemize}
  \item \textbf{\code{operator delete}} --- a class with a \textbf{virtual destructor} names the
    deleting \code{operator delete} in its vtable, even though this code never deletes through a
    base pointer. Every \code{Interface} in the book has one. \textbf{Both} forms are needed: from
    C++14 the compiler calls the sized form but requires the unsized one to exist beside it, and only
    the unsized one is emitted under \code{-fno-sized-deallocation}. On AVR \code{size_t} is
    \code{unsigned int}, so that is the second parameter's type; \code{std::size_t} is not available
    to name it.
  \item \textbf{\code{__cxa_pure_virtual}} --- named by every \textbf{abstract class's} vtable, so
    \code{driver::uart::Interface} alone is enough to require it. It is called only if a pure virtual
    is somehow invoked, so an empty body is honest: there is nothing useful to do and nowhere to
    report it.
  \item \textbf{\code{__cxa_guard_acquire}, \code{_release} and \code{_abort}} --- the guard around a
    \textbf{function-local \code{static}}'s one-time initialization. Single-threaded definitions are
    all a bare-metal target needs.
\end{itemize}

The \code{extern "C"} on the three \code{__cxa_*} groups is not decoration: the compiler emits calls
to them by their C names, and C++ mangling would leave the linker with unresolved symbols.
\code{operator delete} is a C++ operator and must \textbf{not} be \code{extern "C"}. A candidate who
writes empty bodies but mangles the names has not linked anything.

Accept a subset with a note if the candidate says which construct in \emph{their} design each is
for; the mark is for knowing that these are compiler-emitted references with no library behind
them, not for reciting all six.

\textbf{Why it lives in \file{avr/}.} So that the \textbf{host} link keeps libstdc++'s versions.
Replacing the host's thread-safe static guards and its \code{operator delete} with these
single-threaded, hand-written ones would change the behaviour of the build the tests actually run on,
for no benefit whatever. The file exists to satisfy a linker that has no C++ runtime, and the host
has one. \worth{1}

% ==========================================================================================
\answersection{8}{The bench}{12}

\answerpart{a}{4}
\credit{2 marks for the order and completeness of the chain, 1 for the chapters and planes, 1 for
the representations}

\begin{booktable}{@{}rL>{\raggedright\arraybackslash}p{8.4em}>{\raggedright\arraybackslash}p{7.4em}@{}}
  \thead{\#} & \thead{Module} & \thead{Built in} & \thead{Plane} \\
  \midrule
  1 & the USB-serial adapter to the DE0-CV \code{rx} pin & provided hardware & data \\
  2 & \code{sync} on the \code{rx} pin & \chapref{3} & data \\
  3 & \code{uart_rx} & \chapref{3} & data \\
  4 & the RX \code{fifo} & \chapref{4}, instantiated in \chapref{5} & data \\
  5 & \code{uart_regs} & \chapref{5} & the boundary between the two \\
  6 & \code{spi_reg_bridge}, then \code{spi_slave} & provided & control \\
  7 & \code{AvrSpi} & \chapref{9} & control \\
  8 & \code{Uart::readReg}, then \code{Uart::read} & \chapref{7} & control \\
  9 & \code{app::EchoNode::run} & \chapref{10} & neither: it is the application \\
  10 & \code{Uart::write}, then \code{Uart::writeReg} & \chapref{7} & control \\
  11 & \code{AvrSpi}, then \code{spi_slave} and \code{spi_reg_bridge} & \chapref{9}, provided &
    control \\
  12 & \code{uart_regs} and the TX \code{fifo} & Chapters~5 and~4 & the boundary again \\
  13 & the TX feeder in \code{uart_top} & \chapref{5} & data \\
  14 & \code{uart_tx} & \chapref{2} & data \\
  15 & the DE0-CV \code{tx} pin, the adapter, the terminal & provided hardware & data \\
\end{booktable}

\code{baud_gen} (\chapref{2}) paces both 3 and 14 and is on the data plane throughout;
\code{reset_sync} (provided) is on neither and is in every clock domain. \code{uart_top}
(\chapref{1}) is the file that holds all of it together and is not itself a step.

\textbf{The three representations.}
\begin{enumerate}
  \item A \textbf{UART frame on the wire}: a low start bit, eight bits least significant first, a
    high stop bit, timed by the baud rate.
  \item A \textbf{register field over SPI}: the low eight bits of a 32-bit word, carried most
    significant byte first inside a five-byte transaction.
  \item A \textbf{\code{uint8_t}} in the driver and the application.
\end{enumerate}

It changes at \textbf{\code{uart_regs}}, where a frame that has become a FIFO entry becomes a
register field, and again at \textbf{\code{Uart::readReg}}, where four SPI bytes become a
\code{uint32_t} and then a \code{uint8_t}. On the way out it changes back at the same two places, in
the reverse order. Naming those two places is the point of the exercise: everything else on the
list is transport.

\answerpart{b}{3}
\textbf{What the ladder has cleared.}
\begin{itemize}
  \item \textbf{Rung a}, the data-plane pin loopback: the USB-serial adapter, its 3.3\,V logic
    levels, the two data-plane wires and their crossover, and the FPGA pin assignment. None of the
    candidate's logic is in that path, which is what makes it a clean result. What it cannot clear is
    the terminal's baud setting: the terminal talks to itself, and a loopback agrees with itself at
    any rate, for the same reason rung c does.
  \item \textbf{Rung b}, the control plane: the level shifter on all four SPI lines,
    \code{AvrSpi}'s configuration, \code{spi_slave} and \code{spi_reg_bridge}, and the bank's write
    and read paths --- proven by a \code{BAUD_DIV} write and read-back.
  \item \textbf{Rung c}, the peripheral loopback: the whole datapath on real silicon at a real
    50\,MHz --- \code{baud_gen}, \code{uart_tx}, \code{uart_rx}, both FIFOs, the TX feeder and the
    register bank.
\end{itemize}
\worth{1}

\textbf{What is left.} The peripheral talking to a \textbf{different device}, with a \textbf{second,
independently generated baud rate} that now has to agree with the peripheral's own. Rung d is the
first rung at which any absolute timing requirement exists at all.

\textbf{The likely cause, and why rung c hides it.} A \textbf{baud mismatch}. In loopback the
transmitter and receiver are driven by the \emph{same} \code{baud_gen}, from the \emph{same}
50\,MHz clock, with the \emph{same} \code{BAUD_DIV}. Any divider whatsoever passes rung c ---
including a wildly wrong one --- because both ends are wrong by exactly the same amount and the frame
is still sixteen ticks per bit at both ends. Rung d is the first time the number has to be right in
absolute terms.

The two symptoms together point the same way: corruption in both directions is a \emph{timing}
disagreement. A broken, uncrossed or floating data-plane wire gives silence, not wrong characters.
\worth{1}

\textbf{The arithmetic.}

\begin{console}
tick rate = 50_000_000 / 27      = 1_851_851.9 Hz
baud      = 1_851_851.9 / 16     = 115_740.7
error     = (115_740.7 - 115_200) / 115_200 = +0.47%
\end{console}

\noindent Question~3 puts this receiver's combined tolerance at roughly \textbf{+2.5\,\% /
$-$7.7\,\%}, so 0.47\,\% is nowhere near the edge and \code{BAUD_DIV = 27} is \textbf{exonerated}.
To blame the divider you would need a percent-scale error, and the direction matters, because
Question~3 found this receiver's budget lopsided. A \code{BAUD_DIV} of \textbf{28} gives
111\,607~baud, $-3.1$\,\%: this receiver's ticks are then 3.2\,\% longer than the terminal's, beyond
the +2.56\,\% fast-transmitter limit, so typed characters fail while the driver's own bytes, read by
the terminal's receiver with its full five percent, still arrive. A \code{BAUD_DIV} of \textbf{26}
gives 120\,192~baud, +4.3\,\%: inside this receiver's 7.7\,\% on the other side, and only marginal at
the terminal. Neither produces what this bench shows, corruption in \emph{both} directions, which
takes a bigger error: a terminal set to 57600 or 9600 rather than 115200 --- the one setting rung a
cannot check --- or a board clock that is not actually 50\,MHz, which would move both halves of the
peripheral together and so would also have been invisible at rung c. Those are the places to look,
in that order. \worth{1}

\answerpart{c}{3}
\textbf{The \code{main}.}

\begin{cppcode}
/**
 * @brief Application entry point for the ATmega328P.
 */
#include <stdint.h>

#include "app/echo_node.hpp"
#include "driver/transport/avr_spi.hpp"
#include "driver/uart/uart.hpp"

int main()
{
    constexpr uint16_t baudDiv{27U};

    driver::transport::AvrSpi transport{};
    driver::uart::Uart        uart{transport};
    uart.configure(baudDiv);

    app::EchoNode node{uart};

    volatile bool stop{false};
    node.run(stop);

    return 0;
}
\end{cppcode}
\credit{2 marks}

The marks are in the \textbf{order and the ownership}, not the punctuation:
\begin{itemize}
  \item \textbf{Bottom-up construction.} Each layer is built before the one that references it, and
    each is handed up by reference: the transport into \code{Uart}, \code{Uart} into
    \code{EchoNode}. A candidate who constructs \code{EchoNode} first has nothing to give it.
  \item \textbf{\code{configure()} before \code{run()}}, and \code{27} for 115200 --- the one number
    on this page that comes from Part~1's formula rather than from a header. Configure after starting
    the loop and the first bytes go out at whatever the reset divider produces.
  \item \textbf{Nothing is allocated.} Three plain locals in \code{main}'s frame, living until the
    program ends, which on this target is never. This is the freestanding constraint paying off
    rather than biting: the design was reference-injected from \chapref{6}, so the target needs no
    \code{operator new} and none is available. A candidate reaching for \code{new} here has not
    understood what \chapref{9} removed.
  \item \textbf{\code{main} is the only file that names concrete types.} \code{EchoNode} sees an
    interface, \code{Uart} sees an interface; swapping the transport for the \chapref{8} stub is an
    edit to these three lines and nothing else.
\end{itemize}

\code{return 0} is unreachable, and that is fine: \code{run()} never returns with \code{stop} held
\code{false}. Accept an infinite loop after \code{run()} instead, or \code{for(;;)} in place of the
return; do not accept a \code{main} that falls off the end into avr-libc's exit path with the
peripheral still live.

\textbf{Why the flag is \code{false} here.} On the bench there is nobody to stop it: the node should
echo for as long as the board has power, so \code{stop} is a \code{volatile bool} that nothing ever writes
and \code{run()} is an infinite loop by construction. In the host test the very same parameter is
what \emph{ends} the run --- the UART stub of \chapref{6} sets the caller's flag \code{true} the
moment its scripted RX buffer runs out. One flag, taken by \code{const volatile bool&} and read every pass,
serves both: production runs forever because nothing sets it, and the test terminates because
something does. That is why it is a reference rather than a return value or a one-time check.

\textbf{Why not \code{readBlocking()}.} It would spin inside \code{Interface::read()} until a byte
arrived, and would therefore never return to the top of the loop --- so \code{stop} would never be
read again. The application could not be stopped by anybody, including its own test.

That connects directly to how the host test ends: the UART stub of \chapref{6} sets the caller's
\code{stop} flag \code{true} the moment its scripted RX buffer runs out, and \code{run()} sees it on
the \textbf{next pass}. With a blocking read there is no next pass; the test hangs instead of
finishing, which is the test telling you the loop is not actually stoppable. \worth{1}

Worth crediting if a candidate volunteers it: the \emph{write} is allowed to block, and the asymmetry
is deliberate. The byte is already in hand, there is nothing else the loop could usefully do with
it, and waiting cannot deadlock because the TX FIFO is drained by hardware whether software
cooperates or not. Poll where you might have to give up; block where you have already committed.

\answerpart{d}{2}
\textbf{Symptom and risk.} The FPGA input is being driven above its own 3.3\,V supply. At best the
pin misbehaves and the control plane is unreliable --- a \code{BAUD_DIV} write that reads back wrong,
or intermittently right, which is the worse of the two because it looks like a software problem. At
worst the pin's protection diode conducts into the 3.3\,V rail and the pin, or the I/O bank, is
damaged. This is the one genuine electrical hazard on the bench, which is why the wiring is checked
before power is applied.

\textbf{The first rung that exposes it.} \textbf{Rung b}, the control plane. Rung a cannot: the
data-plane pin loopback has no SPI in it at all, running entirely between the adapter and the FPGA,
both at 3.3\,V, with the shifter not in the circuit. That is precisely the property that makes
\textbf{rung a} a useful rung --- it clears the things it clears and nothing else. \worth{1}

\textbf{The safe line.} \textbf{\code{MISO}}. It is the only one of the four that travels from the
FPGA to the Nano, 3.3\,V \emph{up} to 5\,V, so nothing over-drives a 3.3\,V input. It usually works
unshifted because the ATmega's input-high threshold at $V_{cc} = 5$\,V is $0.6 \times V_{cc} =
3.0$\,V, and 3.3\,V clears it.

What can still go wrong: the margin is \textbf{0.3\,V}. A loaded or slow edge, a long breadboard
wire, a supply that has sagged under the FPGA's draw, or a warm room can drop a \code{1} below the
threshold --- producing an intermittent, wiring- and temperature-dependent read failure. That is
considerably worse to debug than a line that never works at all, and it is why ``it mostly works''
is not an acceptable state for this line either. \worth{1}
